\documentclass[11pt]{article}
\usepackage[utf8]{inputenc}
\usepackage{geometry}
\geometry{a4paper, margin=1in}
\usepackage{hyperref}
\usepackage{booktabs}
\usepackage{enumitem}
\usepackage{amsmath}
\usepackage{float}
\usepackage{xcolor}

\title{Lifecycle 1 --- Opening Talk \\ \large How to Start Conveying Your Project to Your Guide}
\author{Speaking Reference}
\date{}

\begin{document}

\maketitle

\noindent \textit{This document is your speaking guide. It tells you what to say, in what order, and why. Read through this before your meeting. The bold sentences are your key talking points --- memorize these.}

\tableofcontents
\newpage

% ============================================================================
\section{Opening (2 minutes) --- The Problem}
% ============================================================================

Start by explaining the real-world problem, not the technology. Your guide wants to know you understand the domain.

\subsection{What to say:}

\textbf{``Landslides kill thousands of people every year and cause billions of dollars in damage. Countries like India, Nepal, China, and the Philippines are most affected because of mountainous terrain and heavy monsoon rainfall.''}

Then add the detection gap:

\textbf{``Currently, landslide identification relies on manual field surveys or expert analysis of satellite images. Both are slow, expensive, and cannot scale. Satellite imagery is widely available --- what is missing is an automated system that can analyze it.''}

\subsection{If your guide asks ``Why is this important?'':}

\begin{itemize}[noitemsep]
    \item India alone loses 500+ lives per year to landslides (NDMA data)
    \item Manual surveys take days to weeks after an event
    \item Early detection could trigger evacuations and save lives
    \item Government agencies like ISRO and NASA already collect satellite imagery --- the bottleneck is analysis, not data
\end{itemize}

% ============================================================================
\section{Why This Project (3 minutes) --- The Gap}
% ============================================================================

Now explain what existing systems do and what they miss. This is the most important part --- it shows your guide that you studied the field before building.

\subsection{What to say:}

\textbf{``I studied existing landslide detection systems and found that they all solve only part of the problem:''}

\begin{enumerate}[noitemsep]
    \item \textbf{Remote sensing systems} (NDVI, change detection) --- detect surface changes but cannot classify what type of landslide occurred or predict future events.
    \item \textbf{ML systems} (SVM, Random Forest on terrain features) --- good for susceptibility mapping but require manual feature engineering and don't work on raw images.
    \item \textbf{Deep learning systems} (U-Net, ResNet for segmentation) --- process images but are single-model, single-phase, and have no temporal prediction component.
\end{enumerate}

Then deliver the key gap:

\textbf{``No existing system combines image-based detection with temporal forecasting. They answer `is this a landslide?' but not `what type is it, how likely is it to recur, and when is the peak risk?'''}

\subsection{If your guide asks ``What papers did you study?'':}

\begin{itemize}[noitemsep]
    \item Meena et al. (2022) --- HR-GLDD dataset paper (our Phase 1 training data)
    \item Lin et al. (2017) --- Focal Loss for Dense Object Detection (our loss function)
    \item Dosovitskiy et al. (2020) --- Vision Transformer (our ViT-B/16 model)
    \item NASA Global Landslide Catalog documentation (our Phase 2 training data)
    \item Tan \& Le (2019) --- EfficientNet: Rethinking Model Scaling (our primary CNN)
\end{itemize}

% ============================================================================
\section{What My Project Does (3 minutes) --- The Solution}
% ============================================================================

\subsection{What to say:}

\textbf{``My project is a two-phase pipeline that solves both detection and prediction:''}

\begin{enumerate}
    \item \textbf{Phase 1 --- Image Classification:} Takes a satellite image and classifies it as landslide or non-landslide using deep learning (CNN). In Lifecycle 1, I started with AlexNet as a baseline and then moved to EfficientNet-B3 for better accuracy with fewer parameters.

    \item \textbf{Phase 2 --- Risk Forecasting:} If Phase 1 detects a landslide, Phase 2 takes the country name as context and uses a Hidden Markov Model trained on 9,471 real NASA landslide events to predict:
    \begin{itemize}[noitemsep]
        \item What type of landslide (debris flow, mudslide, rockfall, etc.)
        \item How likely it is to occur (probability)
        \item When the region is most at risk (peak month)
        \item What types of landslides are likely next (multi-step forecast)
    \end{itemize}
\end{enumerate}

\textbf{``The key contribution is that this is not just a detector --- it is a risk assessment system.''}

% ============================================================================
\section{Why Not Other Projects? (2 minutes) --- Justification}
% ============================================================================

Your guide may ask why you didn't pick a simpler or more common project. Be ready.

\subsection{What to say:}

\textbf{``I chose this project over simpler alternatives for three reasons:''}

\begin{enumerate}
    \item \textbf{Real-world impact:} This is not an academic exercise. Landslide detection is an active research area with government agencies (ISRO, NASA, USGS) investing in it. The results are directly useful.

    \item \textbf{Technical depth:} The project combines two fundamentally different types of machine learning --- deep learning on images (Phase 1) and probabilistic sequential modeling on time series (Phase 2). This makes it significantly more challenging and comprehensive than a standard image classification project.

    \item \textbf{Iterative improvement:} I followed a research methodology --- starting with a simple baseline (AlexNet), measuring its weaknesses, and systematically improving through multiple experiments across three lifecycles. This demonstrates engineering discipline, not just model training.
\end{enumerate}

\subsection{If your guide asks ``Why not just use ResNet/VGG?'':}

\textbf{``I did evaluate them. AlexNet was my starting baseline (62M parameters, simple to explain). I then moved to EfficientNet-B3 which achieves better accuracy with 5$\times$ fewer parameters (12M) due to compound scaling. VGG-16 has 138M parameters and is too slow for practical use. ResNet-50 is good but doesn't scale as efficiently as EfficientNet.''}

\subsection{If your guide asks ``Why not just a classification project?'':}

\textbf{``Pure classification answers only yes/no. Disaster response teams need more --- they need to know what kind of landslide, how dangerous it is, and when it might recur. The HMM phase adds this temporal intelligence using real NASA data spanning 28 years and 141 countries.''}

% ============================================================================
\section{Datasets (2 minutes) --- Why These Datasets}
% ============================================================================

\subsection{What to say:}

\textbf{``I use two publicly available, peer-reviewed datasets:''}

\begin{enumerate}
    \item \textbf{HR-GLDD (Phase 1):} High-Resolution Global Landslide Detection Dataset from Zenodo (DOI: 7189381). Contains satellite image patches at 128$\times$128 pixels with 4 spectral bands (R, G, B, NIR) and pixel-level masks. I converted these to JPEG images for classification --- 2,190 training, 550 validation, 699 test images.

    \item \textbf{NASA Global Landslide Catalog (Phase 2):} 9,471 real landslide events from 1970--2019, covering 141 countries. Each record has the event date, country, landslide type, trigger, size, and fatality count. This is the most comprehensive global landslide event database available.
\end{enumerate}

\subsection{If your guide asks ``Isn't 2,190 images too small?'':}

\textbf{``Yes, which is exactly why I use transfer learning. The models are pre-trained on ImageNet (1.2 million images) which teaches general visual features. I then fine-tune on our 2,190 images which only needs to learn landslide-specific patterns. I also use data augmentation (random flips, rotations, color jitter) and WeightedRandomSampler to handle the 72/28 class imbalance.''}

% ============================================================================
\section{Lifecycle 1 Scope (2 minutes) --- What I Did}
% ============================================================================

\subsection{What to say:}

\textbf{``In Lifecycle 1, I completed three experiments:''}

\begin{enumerate}
    \item \textbf{r0 --- AlexNet from scratch:} Baseline. 78.54\% accuracy, 67.67\% F1-Score, 85.53\% ROC-AUC. This established the performance floor.

    \item \textbf{r1 --- AlexNet with transfer learning:} Added ImageNet pre-trained weights. Accuracy improved to 80.11\% and AUC to 85.73\%. This confirmed that transfer learning helps even with a simple architecture.

    \item \textbf{r2 --- EfficientNet-B3:} Replaced AlexNet with a modern architecture. Recall jumped from 67.30\% to 78.20\% (catches 11\% more landslides), F1 improved to 70.21\%, AUC to 88.93\%. This is the key result of Lifecycle 1.
\end{enumerate}

\subsection{What to say about Lifecycle 2 preview:}

\textbf{``The main gap after Lifecycle 1 is that EfficientNet-B3 is overconfident --- it says 95\% when it is only 75\% correct. In Lifecycle 2, I will fix this with temperature calibration and combine multiple models using an ensemble to further improve accuracy.''}

% ============================================================================
\section{Suggested Speaking Order (Summary)}
% ============================================================================

When you sit with your guide, follow this order:

\begin{enumerate}
    \item \textbf{Problem} (1--2 min): Landslides kill thousands, current detection is manual and slow.
    \item \textbf{Gap} (2--3 min): Existing systems detect but don't classify type or predict future risk.
    \item \textbf{Solution} (2--3 min): Two-phase pipeline --- CNN detection + HMM forecasting.
    \item \textbf{Why this project} (1--2 min): Real-world impact + technical depth + iterative methodology.
    \item \textbf{Datasets} (1--2 min): HR-GLDD for images, NASA GLC for temporal data.
    \item \textbf{LC1 results} (2--3 min): AlexNet baseline $\to$ transfer learning $\to$ EfficientNet-B3. Show metrics table.
    \item \textbf{LC2 preview} (1 min): Temperature calibration + ensemble (already done, just preview it).
    \item \textbf{Demo} (2--3 min): Run a prediction on a test image, show the output.
\end{enumerate}

Total: approximately 15 minutes of speaking.

% ============================================================================
\section{Confidence Builders --- Things You Can Say}
% ============================================================================

These are strong statements you can use if your guide challenges you:

\begin{itemize}
    \item \textbf{``I followed a research methodology, not trial-and-error.''} Each experiment was motivated by measured weaknesses in the previous one.

    \item \textbf{``The datasets are peer-reviewed and publicly available.''} HR-GLDD is from a Zenodo publication, NASA GLC is curated by NASA scientists.

    \item \textbf{``The two-phase architecture is modular.''} Each phase can be improved independently without breaking the other.

    \item \textbf{``I have 7 experiments across 3 lifecycles.''} This shows systematic work, not just running one model.

    \item \textbf{``The HMM adds temporal intelligence that no single-image classifier provides.''} This is the key differentiator from other landslide detection projects.

    \item \textbf{``I used the same training/test split as the original dataset paper.''} This ensures fair and reproducible evaluation.
\end{itemize}

% ============================================================================
\section{Common Mistakes to Avoid}
% ============================================================================

\begin{enumerate}
    \item \textbf{Don't start with technology.} Don't say ``I used EfficientNet-B3 and HMM.'' Start with the problem --- your guide wants to see domain understanding first.

    \item \textbf{Don't memorize numbers blindly.} Know what the numbers mean. If you say ``F1 is 70.21\%'', be ready to explain what F1 measures and why it matters more than accuracy for this task.

    \item \textbf{Don't claim your system is production-ready.} Be honest about limitations --- small dataset, projected r6 metrics, no live satellite integration yet. Honesty builds credibility.

    \item \textbf{Don't compare only accuracy.} When comparing models, always mention Recall and F1. AlexNet pretrained had higher accuracy (80.11\%) than EfficientNet-B3 (79.97\%) but lower recall (67.30\% vs 78.20\%). Accuracy alone is misleading.

    \item \textbf{Don't skip the LC2 preview.} Always end LC1 with a clear statement of what you will do next and why. This shows your work is planned, not ad hoc.
\end{enumerate}

\end{document}
